% !TEX program = xelatex
% Lernplatt - by Can Siebert
% Kurs 71 (Dig 42) - Tag 6 - Aufgabenheft
% Von der Prozesslandkarte zur integrierten Ausfuehrung: BPMS + APIs + Cloud-native

\documentclass[11pt,a4paper,oneside]{article}

\usepackage{polyglossia}
\setdefaultlanguage{german}
\usepackage[a4paper,margin=2.5cm]{geometry}

\usepackage{fontspec}
\defaultfontfeatures{Ligatures=TeX}
\IfFontExistsTF{Charter}{\setmainfont{Charter}}{%
  \IfFontExistsTF{Bitstream Charter}{\setmainfont{Bitstream Charter}}{\setmainfont{TeX Gyre Schola}}%
}
\IfFontExistsTF{Source Sans 3}{\setsansfont{Source Sans 3}}{%
  \IfFontExistsTF{Source Sans Pro}{\setsansfont{Source Sans Pro}}{\setsansfont{TeX Gyre Heros}}%
}
\newcommand{\caveatfontfallback}{\itshape}
\IfFontExistsTF{Caveat}{\newfontfamily\caveatfont{Caveat}[Scale=1.15]}{\newcommand\caveatfont{\caveatfontfallback}}
\setmonofont{Latin Modern Mono}

\usepackage{microtype}
\linespread{1.15}

\usepackage{xcolor}
\definecolor{lpInk}{RGB}{20,28,38}
\definecolor{lpAccent}{RGB}{157,74,26}
\definecolor{lpAccentLight}{RGB}{246,228,212}
\definecolor{lpAmber}{RGB}{222,138,30}
\definecolor{lpAmberLight}{RGB}{255,243,217}
\definecolor{lpAmberBorder}{RGB}{196,118,18}
\definecolor{lpGray}{RGB}{96,104,114}
\definecolor{lpGrayLight}{RGB}{238,240,243}
\definecolor{lpGreen}{RGB}{34,120,72}
\definecolor{lpGreenLight}{RGB}{222,238,228}
\definecolor{lpGreenBorder}{RGB}{24,96,58}
\definecolor{lpL1}{RGB}{34,120,72}
\definecolor{lpL2}{RGB}{22,90,140}
\definecolor{lpL3}{RGB}{170,42,52}

\usepackage[most]{tcolorbox}
\tcbuselibrary{breakable,skins}

\newtcolorbox{infobox}[1][Hinweis]{%
  enhanced, breakable, colback=lpAccentLight, colframe=lpAccent,
  boxrule=0.6pt, arc=2pt, left=10pt,right=10pt,top=8pt,bottom=8pt,
  fonttitle=\sffamily\bfseries\color{white}, coltitle=white,
  title={#1}, attach boxed title to top left={xshift=8pt,yshift=-8pt},
  boxed title style={size=small, colback=lpAccent, colframe=lpAccent, boxrule=0.4pt, arc=1pt},
  top=14pt
}

\newtcolorbox{antwortbox}[1][Ihre Antwort]{%
  enhanced, breakable, colback=white, colframe=lpGray,
  boxrule=0.4pt, arc=2pt, left=10pt,right=10pt,top=8pt,bottom=8pt,
  fonttitle=\sffamily\bfseries\color{white}, coltitle=white,
  title={#1}, attach boxed title to top left={xshift=8pt,yshift=-8pt},
  boxed title style={size=small, colback=lpGray, colframe=lpGray, boxrule=0.4pt, arc=1pt},
  top=14pt
}

\usepackage{enumitem}
\setlist{itemsep=2pt, topsep=4pt, parsep=0pt}
\setlist[itemize,1]{label={\textbullet},leftmargin=1.6em}
\setlist[itemize,2]{label={\textendash},leftmargin=1.4em}
\setlist[enumerate,1]{label=\arabic*., leftmargin=1.8em}

\usepackage{tabularx}
\usepackage{array}
\usepackage{booktabs}
\usepackage{longtable}

\usepackage{fancyhdr}
\usepackage{lastpage}
\pagestyle{fancy}
\fancyhf{}
\renewcommand{\headrulewidth}{0.3pt}
\renewcommand{\footrulewidth}{0pt}
\fancyhead[L]{\footnotesize\sffamily\color{lpGray}Aufgabenheft \textperiodcentered{} Kurs 71 \textperiodcentered{} Tag 6}
\fancyhead[R]{\footnotesize\sffamily\color{lpGray}BPMS \textperiodcentered{} APIs \textperiodcentered{} Cloud-native}
\fancyfoot[C]{\footnotesize\sffamily\color{lpGray}\thepage{} / \pageref{LastPage}}

\fancypagestyle{plain}{%
  \fancyhf{}%
  \renewcommand{\headrulewidth}{0pt}%
  \fancyfoot[C]{\footnotesize\sffamily\color{lpGray}\thepage{} / \pageref{LastPage}}%
}

\usepackage[explicit]{titlesec}
\titleformat{\section}[block]
  {\sffamily\Large\bfseries\color{lpAccent}}
  {\thesection.}{0.6em}{#1}
  [{\vspace{2pt}\color{lpAccent}\titlerule[0.6pt]}]
\titleformat{\subsection}[block]
  {\sffamily\large\bfseries\color{lpInk}}
  {\thesubsection}{0.6em}{#1}
\titlespacing*{\section}{0pt}{14pt}{8pt}
\titlespacing*{\subsection}{0pt}{10pt}{4pt}

\usepackage[hidelinks,unicode]{hyperref}

\newcommand{\akzent}[1]{{\caveatfont\color{lpAmber}#1}}
\newcommand{\fachbegriff}[1]{\textbf{#1}}
\newcommand{\quelle}[1]{{\footnotesize\sffamily\color{lpGray}(#1)}}

\newcommand{\niveaubadge}[2]{%
  \tcbox[on line, colback=#1, colframe=#1, coltext=white,
         boxrule=0pt, arc=2pt, left=5pt,right=5pt,top=1pt,bottom=1pt,
         nobeforeafter]{\sffamily\bfseries\footnotesize #2}%
}
\newcommand{\nivLi}{\niveaubadge{lpL1}{L1\,\textbullet\,Wissen}}
\newcommand{\nivLii}{\niveaubadge{lpL2}{L2\,\textbullet\,Anwendung}}
\newcommand{\nivLiii}{\niveaubadge{lpL3}{L3\,\textbullet\,Management}}

\newcommand{\zeit}[1]{%
  \tcbox[on line, colback=lpGrayLight, colframe=lpGrayLight, coltext=lpInk,
         boxrule=0pt, arc=2pt, left=5pt,right=5pt,top=1pt,bottom=1pt,
         nobeforeafter]{\sffamily\footnotesize\textbf{$\sim$\,#1}}%
}

\newcommand{\aufgabenkopf}[4]{%
  \par\vspace{6pt}
  \noindent{\sffamily\Large\bfseries\color{lpAccent}Aufgabe #1 \textendash{} #2}\par
  \vspace{2pt}
  \noindent{\color{lpAccent}\rule{\linewidth}{0.6pt}}\par
  \vspace{4pt}
  \noindent #3 \hfill \zeit{#4}\par
  \vspace{6pt}
}

\newcommand{\ytquery}[1]{%
  \par\vspace{4pt}
  \noindent{\footnotesize\sffamily\color{lpGray}\textbf{YouTube-Suchquery:} \texttt{#1}}\par
}

\newcommand{\antwortlinien}[1]{%
  \par\vspace{6pt}
  \foreach \x in {1,...,#1}{%
    \noindent\makebox[\linewidth]{\dotfill}\\[6pt]
  }
}
\usepackage{pgffor}

\begin{document}

\begin{titlepage}
  \pagestyle{empty}
  \vspace*{1.5cm}
  \noindent{\sffamily\footnotesize\color{lpGray}LERNPLATT \textperiodcentered{} BY CAN SIEBERT}\\[2pt]
  \noindent{\color{lpAccent}\rule{\linewidth}{1.2pt}}\\[4pt]
  \noindent{\sffamily\footnotesize\color{lpGray}AUFGABENHEFT \textperiodcentered{} MODUL 2 \textperiodcentered{} TAG 6 VON 16 \textperiodcentered{} KURS 71 (Dig 42) \textperiodcentered{} 5.\,Juni 2026}

  \vspace{3cm}
  {\sffamily\fontsize{30}{34}\selectfont\bfseries\color{lpInk}Aufgabenheft\\Tag 6\par}

  \vspace{0.7cm}
  {\sffamily\large\color{lpGray}Von der Prozesslandkarte zur\\integrierten Ausf\"uhrung \textendash{} BPMS,\\APIs und Cloud-native Skalierung.\par}

  \vspace{1.5cm}
  {\caveatfont\LARGE\color{lpAmber}11 Aufgaben \textperiodcentered{} L1 \textperiodcentered{} L2 \textperiodcentered{} L3}\par

  \vfill

  \noindent\begin{minipage}{0.55\linewidth}
    {\sffamily\footnotesize\color{lpGray}Niveau-Mix:}\\
    {\sffamily\small\color{lpInk}3\,\textsf{L1} \textperiodcentered{} 5\,\textsf{L2} \textperiodcentered{} 3\,\textsf{L3}}\\[10pt]
    {\sffamily\footnotesize\color{lpGray}Bearbeitung:}\\
    {\sffamily\small\color{lpInk}Selbstbearbeitung im Lernpfad}
  \end{minipage}\hfill
  \begin{minipage}{0.4\linewidth}
    \raggedleft
    {\sffamily\footnotesize\color{lpGray}Dozent:}\\
    {\sffamily\small\color{lpInk}Can Siebert}\\[10pt]
    {\sffamily\footnotesize\color{lpGray}Begleitend:}\\
    {\sffamily\small\color{lpInk}Lehrskript \& L\"osungsheft}
  \end{minipage}

  \vspace{0.5cm}
  \noindent{\color{lpAccent}\rule{\linewidth}{0.6pt}}
\end{titlepage}

\section*{So arbeiten Sie mit diesem Heft}
\thispagestyle{plain}

\noindent Dieses Aufgabenheft enth\"alt elf Aufgaben zu Tag 6 \textendash{} \emph{Von der Prozesslandkarte zur integrierten Ausf\"uhrung}: BPMS, APIs, Cloud-native Skalierung und E2E-Methodenintegration. Das Heft baut auf den Repository- und Execution-Themen von Tag 4 auf und erweitert sie um Integration, Plattform-Skalierung und Operating-Model-Fragen. Die Aufgaben gliedern sich in drei Niveaus:

\begin{itemize}
  \item \nivLi{} \textendash{} \fachbegriff{Wissen.} BPMS-Bausteine, API-Vertrag, Cloud-native Prinzipien. 5\,--\,10 Minuten.
  \item \nivLii{} \textendash{} \fachbegriff{Anwendung.} BPMS-Blueprint, API-Vertrag, Sync/Async-Entscheidung. 15\,--\,30 Minuten.
  \item \nivLiii{} \textendash{} \fachbegriff{Management.} Fallstudie \emph{EuroLogix}, Plattform-Strategie. 30\,--\,45 Minuten.
\end{itemize}

\noindent Jede Aufgabe enth\"alt:
\begin{itemize}
  \item eine Niveau-Markierung und einen Zeit-Richtwert,
  \item den Aufgabentext (oft mit Kontext oder Fallbeschreibung),
  \item einen Antwort-Freiraum f\"ur Ihre L\"osung,
  \item eine \emph{YouTube-Suchquery} f\"ur den begleitenden Lernpfad.
\end{itemize}

\noindent\textbf{Architekturentscheidungs-Notiz (ADR):} Halten Sie wichtige Entscheidungen kurz strukturiert fest: \emph{Kontext}, \emph{Entscheidung}, \emph{Begr\"undung}, \emph{verworfene Alternativen}, \emph{Fr\"uhwarnsignale} f\"ur Re-Evaluation. Eine ADR ist nicht \"uberbordend \textendash{} eine halbe A4-Seite reicht.

\begin{infobox}[Hinweis]
Die Musterl\"osungen finden Sie im \emph{L\"osungsheft Tag 6}. Heute steht die Frage im Mittelpunkt: \emph{Wer tr\"agt die Verantwortung}, wenn Prozesse \"uber Plattformen laufen? BPMS, APIs und Cloud-native sind Plattform-Entscheidungen \textendash{} keine Tool-Auswahl. Insbesondere die L3-Aufgaben verlangen eine Architekturentscheidung, die in 6\,Monaten teuer zu \"andern ist \textendash{} dokumentieren Sie deshalb Annahmen, verworfene Alternativen und Fr\"uhwarnsignale (\emph{ADR-Format}).
\end{infobox}

\clearpage

\section*{Niveau L1 \textendash{} Wissen \& Reproduktion}
\addcontentsline{toc}{section}{L1 -- Wissen}

% LERNPLATT_HILFSVIDEOS
\section*{Hilfsvideos zu diesem Tag}
\addcontentsline{toc}{section}{Hilfsvideos zu diesem Tag}
\thispagestyle{plain}

\noindent Die folgenden YouTube-Erkl\"arvideos vermitteln die Inhalte, die Sie zur Bearbeitung der Aufgaben ben\"otigen. Klicken Sie auf den Titel, um das Video direkt zu \"offnen; die vollst\"andige URL steht in Klammern.

\begin{description}[style=nextline,leftmargin=18pt,labelindent=0pt,itemsep=8pt]
  \item[1.~Was ist BPM? — Business Process Management in 2 Minuten] \href{https://youtu.be/9NXq0a-X7cI}{\textcolor{lpAccent}{https://youtu.be/9NXq0a-X7cI}}\\[2pt]
  {\footnotesize\color{lpGray}(URL: \nolinkurl{https://youtu.be/9NXq0a-X7cI})}
  \item[2.~Was ist eine REST-API? — Einführung in REST APIs] \href{https://youtu.be/ZuINS_-n-AM}{\textcolor{lpAccent}{https://youtu.be/ZuINS_-n-AM}}\\[2pt]
  {\footnotesize\color{lpGray}(URL: \nolinkurl{https://youtu.be/ZuINS_-n-AM})}
  \item[3.~Monolithen vs. Microservices — wann ist welche Architektur sinnvoll?] \href{https://youtu.be/N0k3VKL4Als}{\textcolor{lpAccent}{https://youtu.be/N0k3VKL4Als}}\\[2pt]
  {\footnotesize\color{lpGray}(URL: \nolinkurl{https://youtu.be/N0k3VKL4Als})}
  \item[4.~Business Process Management — Closed-Loop einfach erklärt] \href{https://youtu.be/IBenKthHmUU}{\textcolor{lpAccent}{https://youtu.be/IBenKthHmUU}}\\[2pt]
  {\footnotesize\color{lpGray}(URL: \nolinkurl{https://youtu.be/IBenKthHmUU})}
\end{description}
\vspace{0.6em}
\noindent{\footnotesize\color{lpGray}\textit{Tipp:} \"Offnen Sie das Video, lassen Sie es im Hintergrund laufen und arbeiten Sie parallel an der Aufgabe.}

\clearpage

\aufgabenkopf{1}{BPMS-Bausteine}{\nivLi}{8\,min}

\noindent Ein \fachbegriff{BPMS} (Business Process Management System) ist mehr als eine Engine. Listen Sie die \emph{sieben Kernbausteine} und beschreiben Sie pro Baustein: \emph{Zweck} und \emph{Risiko bei Fehlen}.

\medskip
\noindent\textbf{Hinweis:} \emph{Prozessengine}, \emph{Aufgaben-/Worklist}, \emph{Regelwerk (DMN)}, \emph{Forms}, \emph{Monitoring/Dashboards}, \emph{Audit Trail}, \emph{Versionierung}.

\medskip
\noindent Erg\"anzen Sie: Was bedeutet ``Closed Loop'' im BPMS-Kontext? Beschreiben Sie die vier Phasen \emph{Modellieren \textrightarrow{} Ausf\"uhren \textrightarrow{} \"Uberwachen \textrightarrow{} Optimieren}.

\ytquery{BPMS Bausteine Prozessengine Workflow Engine Closed Loop Camunda}

\antwortlinien{14}

\clearpage

\aufgabenkopf{2}{API als Vertrag}{\nivLi}{10\,min}

\noindent Ein API-Aufruf ist nicht ein ``Quick Call'', sondern ein \fachbegriff{Vertrag}. Listen Sie die acht Pflichtbestandteile eines API-Vertrags:

\medskip
\noindent\textbf{Hinweis:} \emph{Endpoint}, \emph{Input}, \emph{Output}, \emph{Fehlercodes}, \emph{Auth}, \emph{Rate Limits}, \emph{Versionierung}, \emph{Monitoring / Correlation-ID}.

\medskip
\noindent F\"ur jeden Bestandteil:
\begin{enumerate}
  \item Was wird vereinbart?
  \item Was passiert, wenn er \emph{fehlt} oder \emph{schwammig} formuliert ist?
\end{enumerate}

\noindent Schlie{\ss}en Sie mit einem Satz: Wer ist Eigent\"umer des API-Vertrags \textendash{} \emph{Anbieter} oder \emph{Konsument}?

\ytquery{API Contract REST Endpoint Versioning Auth Rate Limiting Best Practice}

\antwortlinien{14}

\clearpage

\aufgabenkopf{3}{Sync vs.\ Async, Idempotenz, Retry}{\nivLi}{8\,min}

\noindent Erkl\"aren Sie pr\"agnant (je 2\,--\,3 S\"atze):

\begin{enumerate}
  \item \emph{Sync API Call} vs.\ \emph{Async Event/Message}: Wann was?
  \item \emph{Idempotenz}: Was hei{\ss}t es? Warum ist sie f\"ur Retry-F\"ahigkeit zwingend?
  \item \emph{Retry-Strategie}: Was bedeutet ``exponential backoff''? Wann darf man \emph{nicht} retryen (Side-Effects)?
  \item \emph{Correlation / Case ID}: Welche Funktion erf\"ullt sie im Async-Modus?
\end{enumerate}

\medskip
\noindent Erg\"anzen Sie eine Vergleichstabelle Sync vs.\ Async:

\medskip
\begin{tabularx}{\linewidth}{@{}p{3.6cm}|X|X@{}}
\toprule
\textbf{Aspekt} & \textbf{Sync} & \textbf{Async} \\
\midrule
Antwortverhalten & & \\[14pt]
\midrule
Kopplung         & & \\[14pt]
\midrule
Resilienz        & & \\[14pt]
\midrule
Komplexit\"at    & & \\[14pt]
\bottomrule
\end{tabularx}

\ytquery{Sync vs Async API Event Driven Architecture Idempotency Retry Exponential}

\antwortlinien{6}

\clearpage

\section*{Niveau L2 \textendash{} Anwendung \& Transfer}
\addcontentsline{toc}{section}{L2 -- Anwendung}

\aufgabenkopf{4}{BPMS-Blueprint \emph{Customer Complaint Handling}}{\nivLii}{25\,min}

\begin{infobox}[Textprozess]
\emph{``Customer Complaint Handling''}: Reklamation kommt rein \textrightarrow{} Klassifizieren \textrightarrow{} ggf.\ R\"uckfrage \textrightarrow{} Entscheidung Kulanz/Retoure \textrightarrow{} Abwicklung \textrightarrow{} Abschluss \textrightarrow{} Reporting. SLAs: 48\,h Erstreaktion, 10\,Tage Abschluss. Daten in CRM, ERP, E-Mail.
\end{infobox}

\noindent\textbf{Arbeitsauftrag:}

\begin{enumerate}
  \item \textbf{BPMS-Blueprint-Skizze:}
    \begin{itemize}
      \item \emph{Human Tasks} (wer?),
      \item \emph{Systemaktionen} (welches System?),
      \item \emph{Regeln} (welche Entscheidung?),
      \item \emph{SLAs / Eskalation} (wann?).
    \end{itemize}
  \item \textbf{3 KPIs} (z.\,B.\ SLA-Einhaltung, Rework-Quote, Kundenzufriedenheit) + Owner.
  \item \textbf{Entscheidung unter Unsicherheit:} Welche \emph{eine Funktion} \emph{muss} im BPMS sein, welche darf zun\"achst au{\ss}erhalb bleiben (und warum)?
\end{enumerate}

\ytquery{BPMS Blueprint Customer Complaint Workflow SLA Eskalation Camunda}

\antwortlinien{20}

\clearpage

\aufgabenkopf{5}{API-Vertr\"age f\"ur den Complaint-Prozess}{\nivLii}{22\,min}

\begin{infobox}[Integrationsbedarf]
Der Complaint-Prozess braucht Integrationen zu \emph{CRM} (Kunde), \emph{ERP} (Auftrag/Rechnung), \emph{Versand} (Tracking).
\end{infobox}

\noindent\textbf{Arbeitsauftrag:}

\begin{enumerate}
  \item Definieren Sie \emph{zwei APIs}, die Sie zwingend brauchen (z.\,B.\ \texttt{GetOrderStatus}, \texttt{CreateReturn}).
  \item F\"ullen Sie pro API den Vertrag aus: \emph{Name/Zweck}, \emph{Input-Felder} (mind.\ 3), \emph{Output-Felder} (mind.\ 2), \emph{Fehlerf\"alle} (mind.\ 2), \emph{Auth}, \emph{Versionierung}, \emph{Monitoring / Correlation-ID}.
  \item Definieren Sie \emph{3 Fehlerf\"alle} (Timeout, Daten fehlen, Auth-Fail) und wie der Prozess reagiert (\emph{Retry}, \emph{manuell}, \emph{Eskalation}).
  \item \textbf{Sync- oder Async-Entscheidung} pro API: mit Begr\"undung (Speed vs.\ Robustheit).
\end{enumerate}

\ytquery{REST API Contract Order Status Return GetOrderStatus CreateReturn Error}

\antwortlinien{18}

\clearpage

\aufgabenkopf{6}{Cloud-native Betriebsanforderungen}{\nivLii}{20\,min}

\begin{infobox}[Anforderungen Logistik]
Das Volumen schwankt stark (Peak-Saisons), das System ist auditiert, Kunden erwarten Real-time-Tracking.
\end{infobox}

\noindent\textbf{Arbeitsauftrag:}

\noindent Leiten Sie \emph{f\"unf cloud-native Betriebsanforderungen} ab, mit je einer 1-Satz-Begr\"undung:

\begin{enumerate}
  \item \emph{Auto-Scaling} (horizontal) \textendash{} warum, woran erkennbar?
  \item \emph{Observability} (Logs / Metrics / Traces) \textendash{} welche drei Daten zwingend?
  \item \emph{Resilienz} (Circuit Breaker, Retry) \textendash{} welcher Schutz vor Kaskaden?
  \item \emph{Deployment} (CI/CD, Blue/Green) \textendash{} wie minimieren Sie Release-Risiko?
  \item \emph{Security} (Zero Trust konzeptionell) \textendash{} welche zwei Prinzipien tragen?
\end{enumerate}

\noindent \emph{Anti-Muster}: Skalierung ohne Observability ist \emph{Blindflug}. Begr\"unden Sie in 3 S\"atzen.

\ytquery{Cloud Native Twelve Factor Auto Scaling Observability Circuit Breaker SLO}

\antwortlinien{18}

\clearpage

\aufgabenkopf{7}{System-of-Record-Entscheidung}{\nivLii}{18\,min}

\begin{infobox}[Mini-Szenario]
F\"ur den Complaint-Prozess gibt es f\"ur das Datenobjekt \emph{Kunde} potenziell mehrere Master: \emph{CRM} (vertriebsseitig), \emph{ERP} (rechnungsseitig), \emph{Webshop} (selbstgepflegt).
\end{infobox}

\noindent\textbf{Arbeitsauftrag:}

\begin{enumerate}
  \item Definieren Sie das \fachbegriff{System of Record (SoR)} pro Datenobjekt:
    \begin{itemize}
      \item \emph{Kunde} \textendash{} welches System ist Master?
      \item \emph{Auftrag} \textendash{} welches System?
      \item \emph{Reklamation} \textendash{} welches System (BPMS? Ticketing?)?
    \end{itemize}
  \item F\"ur \emph{ein} Datenobjekt: Was passiert bei Konflikten zwischen zwei Systemen?
  \item Definieren Sie eine \emph{Master-Slave-Regel} oder eine Synchronisationsstrategie.
  \item Entscheidung unter Unsicherheit: Welche \emph{eine Inkonsistenz} akzeptieren Sie bewusst tempor\"ar (mit Guardrail) \textendash{} und welche nie?
\end{enumerate}

\ytquery{System of Record Data Ownership Master Data Management CRM ERP Conflict}

\antwortlinien{14}

\clearpage

\aufgabenkopf{8}{E2E-Methodenintegration: Model \textrightarrow{} Run \textrightarrow{} Improve}{\nivLii}{18\,min}

\begin{infobox}[Idee]
E2E hei{\ss}t: Prozess + Daten + Systeme + Rollen + KPIs + Controls werden zusammengef\"uhrt. \emph{``Model \textrightarrow{} Run \textrightarrow{} Improve''} ist die Steuerungsschleife dahinter.
\end{infobox}

\noindent\textbf{Arbeitsauftrag:}

\begin{enumerate}
  \item \emph{Model:} Welche zwei Modellierungs-Sichten brauchen Sie f\"ur den Complaint-Prozess (BPMN? VSM? Beides)?
  \item \emph{Run:} Welche BPMS-Komponente plus welche API-Integration bringt die Modelle zum Laufen?
  \item \emph{Improve:} Welche Mining-/KPI-Quelle schlie{\ss}t den Loop? Konkrete Idee f\"ur den ersten Review-Zyklus.
  \item \emph{Process Review Board:} Wer sitzt am Tisch (Rollen) und wie oft (Frequenz)?
\end{enumerate}

\noindent Faustregel (1 Satz): Woran erkennen Sie, dass der Loop ``geschlossen'' ist \textendash{} im Gegensatz zu drei isolierten Aktivit\"aten?

\ytquery{Process Model Run Improve Closed Loop BPMS Mining KPI Continuous Improvement}

\antwortlinien{14}

\clearpage

\section*{Niveau L3 \textendash{} Management \& Entscheidungsarchitektur}
\addcontentsline{toc}{section}{L3 -- Management}

\aufgabenkopf{9}{Fallstudie \emph{EuroLogix}}{\nivLiii}{45\,min}

\begin{infobox}[Fallstudie \emph{EuroLogix}]
\textbf{Unternehmen:} \emph{EuroLogix} (Logistikdienstleister).\\
\textbf{Problem:} E2E \emph{Shipment Exception Handling} ist langsam, viele Medienbr\"uche, Kundenbeschwerden. Systeme: TMS, CRM, E-Mail, Data Warehouse. Volumen schwankt stark.

\smallskip
\textbf{Restriktionen:} \textbf{9 Wochen}, Budget \textbf{150\,k\,\euro}, Legacy-Systeme, Security-Anforderungen, unklare Peak-Lastdaten.
\end{infobox}

\noindent\textbf{Arbeitsauftrag:}

\begin{enumerate}
  \item \textbf{BPMS-gest\"utzter Prozess} inkl.\ Rollen, SLAs, Eskalationen, Audit Trail (mind.\ 2 Eskalationsregeln).
  \item \textbf{3 Integrationspunkte} (APIs/Events) inkl.\ Vertragsfeldern und Fehler-/Fallback-Strategie.
  \item \textbf{Cloud-native Betriebsanforderungen} (mind.\ 4 Punkte): Skalierung, Observability, Resilienz, Deployment.
  \item \textbf{E2E-Methodenkonzept:} BPMN-Soll + VSM-Engpass + KPIs + Mining \textrightarrow{} Process Review Board.
  \item \textbf{MVP-Entscheidung} (Release 1 in 9 Wochen): was ist drin, was \emph{explizit nicht}?
\end{enumerate}

\noindent\textbf{Format-Hinweis:} eine A4-Seite Senior-Pitch. In 90\,Sekunden lesbar.

\ytquery{Logistics Exception Handling BPMS Cloud Native API Integration Roadmap}

\antwortlinien{36}

\clearpage

\aufgabenkopf{10}{Architekturpfad: event-driven vs.\ sync}{\nivLiii}{25\,min}

\noindent\textbf{Diskussionsaufgabe.}

\noindent Eine fundamentale Architekturentscheidung: \emph{\"uberwiegend event-driven} (Choreographie, lose Kopplung) oder \emph{\"uberwiegend sync} (Orchestrierung, klare Reihenfolge). Beide haben Folgen f\"ur 6+ Monate.

\medskip
\noindent\textbf{Arbeitsauftrag:}

\begin{enumerate}
  \item \textbf{Trade-off-Matrix} entlang 5 Dimensionen: \emph{Resilienz}, \emph{Nachvollziehbarkeit / Audit}, \emph{Skalierung}, \emph{Komplexit\"at}, \emph{Team-Skills}.
  \item \textbf{Entscheidung} f\"ur \emph{EuroLogix} mit Begr\"undung in 4\,--\,5 S\"atzen.
  \item \textbf{Hybrid:} Welche zwei Use Cases bleiben bewusst \emph{sync} (auch wenn der Rest event-driven ist)?
  \item \textbf{Fr\"uhwarnsignal:} Was zwingt Sie in 6 Monaten zur Umkehr?
\end{enumerate}

\ytquery{Event Driven Architecture vs Synchronous Trade Off Resilience Audit}

\antwortlinien{18}

\clearpage

\aufgabenkopf{11}{Transferprojekt \textendash{} E2E Integration \& BPMS Platform Strategy}{\nivLiii}{45\,min}

\begin{infobox}[Rolle und Ausgangslage]
\textbf{Ihre Rolle:} Chief Digital Officer / Enterprise Architect / Head of Platform Engineering.

\smallskip
\textbf{Auftrag:} In \textbf{16 Wochen} ein skalierbares E2E-Integrations- und Ausf\"uhrungsmodell etablieren (BPMS + API + Cloud-native), das Prozesssteuerung messbar macht und Audit-/Security-Anforderungen erf\"ullt.

\smallskip
\textbf{Restriktionen:} Legacy-Landschaft, Peak-Lasten unbekannt, Budget \textbf{250\,k\,\euro}, Team-Kapazit\"aten schwanken, regulatorischer Druck.
\end{infobox}

\noindent\textbf{Arbeitsauftrag:}

\begin{enumerate}
  \item \textbf{Plattform-Entscheidungsarchitektur:} Welche Entscheidungen werden \emph{zentral} (Standards) vs.\ \emph{dezentral} (Teams) getroffen? (APIs, Prozessstandards, Observability, Security.)
  \item \textbf{API Governance:} Versionierung, SLA/SLO, Ownership, Contract Testing (konzeptionell), Change-Prozess.
  \item \textbf{BPMS Governance:} Prozess-Lifecycle, Release-Management, Exception-Policy, Audit Trail, KPI-Steuerung.
  \item \textbf{Cloud-native Betriebsmodell:} SLOs, Skalierungs\-strategie, Incident Management, Kostenkontrolle (FinOps light).
  \item \textbf{Roadmap \& Investment Cases:} 3 Releases, Nutzen-/Risiko\-argumentation, Messkonzept (Leading / Lagging).
  \item \textbf{Zwei Entscheidungen unter Unsicherheit:}
    \begin{itemize}
      \item \emph{Architekturpfad:} event-driven vs.\ \"uberwiegend sync \textendash{} Entscheidung + Risiko + Fr\"uhwarnsignal.
      \item \emph{MVP-Scope:} welche \emph{2 Integrationen} zuerst \textendash{} Entscheidung + Begr\"undung + ``Stop-Doing''-Liste.
    \end{itemize}
\end{enumerate}

\noindent\textbf{Bewertungskriterien:} Strategie-Anschluss \textperiodcentered{} Klarheit der Verantwortung \textperiodcentered{} Pragmatismus unter Restriktionen \textperiodcentered{} Entscheidungslogik.

\ytquery{Platform Strategy BPMS Integration Cloud Native API Governance CDO}

\antwortlinien{38}

\clearpage

\section*{Abschluss}
\thispagestyle{plain}

\noindent Sie haben alle elf Aufgaben bearbeitet. Vergleichen Sie nun Ihre Antworten mit dem \emph{L\"osungsheft Tag 6}. Pr\"ufen Sie:

\begin{itemize}
  \item Habe ich die BPMS-Bausteine als \emph{Closed Loop} verstanden oder als Tool-Auflistung?
  \item Sind meine API-Vertr\"age \emph{vollst\"andig} (Versionierung, Auth, Monitoring) \textendash{} oder fehlen die Compliance-Bestandteile?
  \item Habe ich Sync und Async \emph{bewusst} entschieden \textendash{} oder per Default ``sync, weil einfacher''?
  \item Habe ich Cloud-native nicht nur als ``Hosting in der Cloud'' verstanden, sondern als Betriebsprinzip (Stateless, Observability, Auto-Scaling)?
  \item Hat mein System-of-Record klare Konflikt-Regeln \textendash{} oder bleibt es ``klingt schon ok''?
\end{itemize}

\medskip
\noindent\textbf{Senior-Frage:} Welche \emph{eine} Entscheidung heute w\"urden Sie auch dann verteidigen, wenn das Audit in 6 Monaten unangek\"undigt kommt? Genau diese Entscheidung dokumentieren Sie heute als \emph{Architectural Decision Record (ADR)} mit Annahmen, verworfener Alternative und Fr\"uhwarnsignal.

\medskip
\begin{infobox}[Was Sie jetzt k\"onnen sollten]
\begin{itemize}
  \item BPMS-Bausteine als geschlossenes Steuerungssystem (Closed Loop) einsetzen.
  \item API-Vertr\"age vollst\"andig formulieren \textendash{} mit Versionierung, Auth, Monitoring, Fehlerf\"allen.
  \item Sync- vs.\ Async-Entscheidungen bewusst und mit Begr\"undung treffen.
  \item Cloud-native Betriebsanforderungen (Skalierung, Observability, Resilienz) ableiten.
  \item System of Record pro Datenobjekt definieren \textendash{} und Konflikte ausregeln.
  \item Eine Plattform-Strategie schreiben, die zentralen Standards mit dezentralen Teams verbindet.
\end{itemize}
\end{infobox}

\vspace{1cm}
\noindent{\color{lpAccent}\rule{\linewidth}{0.6pt}}\\[4pt]
{\footnotesize\sffamily\color{lpGray}Lernplatt \textperiodcentered{} by Can Siebert \textperiodcentered{} Kurs 71 (Dig 42) \textperiodcentered{} Tag 6 \textperiodcentered{} Aufgabenheft \textperiodcentered{} \textcopyright{} 2026}

\end{document}
